\section{A price with an extra right}
\label{sec:deamright}

Chapter~6 ended with a debt.  The parity line on the running node ---
SPY December 2026, the (underlier, expiry) pair every chapter has
worked on --- was almost perfectly straight: its residuals ran to
\MacFwdLineStraightPct\% of the full range the observable spans.
But those residuals were not noise.  They bowed, cresting at
$+\$\MacFwdLineResidMaxDollars$ near the money, hundreds of times the
one-cent tick (\cref{fig:fwdline}b).  Chapter~6 identified the culprit
in one sentence --- SPY's listed options are American --- and promised
that this chapter would measure and remove it.  We now pay the debt.

Start with the vocabulary, because everything follows from it.  A
\emph{European} option can be exercised only at expiry; every model in
Part~I priced European options, as terminal expectations under one
probability law (Chapter~2).  An \emph{American} option can be
exercised at any moment its holder chooses, up to and including
expiry.  US-listed single-stock and ETF options --- the raw quotes
$\mathsf{C}(K)$, $\mathsf{P}(K)$ this book's frozen chains carry ---
are American.  Index options, and the synthetic test boards the book
constructs for itself, are European; for them, every step in this
chapter leaves the quotes untouched.

The extra right cannot make the option cheaper.  Whatever plan a
European holder has --- namely, wait to expiry --- an American holder
may copy; the American holder simply has more plans to choose from,
and the best of a larger set of plans is worth at least the best of a
smaller one.  Write $A(\sigma)$ for the model price of one American
option as a function of the volatility $\sigma$, everything else (its
strike, expiry, and carry) held fixed, and $E(\sigma)$ for its
\emph{European twin}: the same contract with the early-exercise right
removed.  Then
\[
A(\sigma)\;\ge\;E(\sigma),
\]
and the difference $A-E$ --- the \emph{early-exercise premium}, the
dollar value of the right itself --- is never negative.  It is often
exactly zero, and \cref{sec:deamwhere} proves where.  By the same
larger-set logic, an American price can never sit below the option's
\emph{intrinsic value} --- what exercising this instant would pay,
$(S-K)^+$ for a call and $(K-S)^+$ for a put --- because ``exercise
this instant'' is itself one of the holder's plans.
\Cref{sec:deamwhere} makes both statements precise.  One notation
note, once: italic $E$ is always a \emph{price} in this book; the
expectation operator keeps its blackboard $\E$.  The premium itself
gets no symbol --- we write $A-E$, or say ``the premium''.

Here is why the premium matters to a volatility fitter.  The market
prints one number per quote: the American price.  It does not print
$E$, and it does not print the premium; the decomposition is real but
unobservable, quote by quote.  Part~I's machinery inverts a
\emph{European} price into an implied volatility --- that is what the
Black formula's argument is --- so feeding it an American mid hands it
a price that is too high for what it thinks it is pricing.  The
inversion has no slot for the premium, so it books the whole excess as
volatility, and the smile comes out biased upward.

How badly?  Recall the dictionary from Chapter~6: a price error
$\Delta P$ moves an inverted volatility by roughly $\Delta
P/\text{vega}$, where \emph{vega} --- the price sensitivity to
volatility, $\partial C/\partial\sigma = DF\,\phin(d_+)\sqrt{\tau}$
for the dollar call, with $d_\pm$ the Black formula's usual arguments
(first courses write $d_1$, $d_2$) --- peaks at the money and
collapses both deep in the money and far out of it.  A premium of fixed dollar size therefore
reads as a \emph{large} phantom volatility exactly where vega is
small: deep in the money.  The damage is measured in vol basis points
(one vol bp is $10^{-4}$ of absolute volatility, as in Chapter~6) not
because the dollars are large --- they are cents to a few dollars ---
but because the wings divide by a small vega.

\Cref{fig:deamwedge} runs the whole experiment on a synthetic board
where the truth is known: American calls on a stock paying a $6\%$
dividend yield against a $2\%$ rate, half a year to expiry, priced by
this chapter's tree at a true volatility of exactly $25\%$ (all
constants in appendix~\ref{app:deamprotocol}; why dividends make call
premiums live is \cref{sec:deamwhere}'s subject).  In panel~(a), the
orange curve is the naive reading: each American price handed to the
analytic European inversion.  At the forward the phantom is already
\MacDeamWedgeAtmBiasBp\ vol bp; moving into the money it grows without
mercy, reaching \MacDeamWedgeMaxBiasBp\ vol bp at the deepest strike.
The blue curve is the subtraction this chapter builds
(\cref{sec:deamsub}): it returns a flat $25\%$ across the board, to an
rms of \MacDeamWedgeRootRmsBp\ vol bp.  Notice where the blue curve
\emph{stops}: below \$\MacDeamMapCallBdryPct\ there are no blue points
at all, because
there the price has collapsed onto its exercise floor and contains no
volatility to recover --- the dead zone \cref{sec:deamwhere} maps and
\cref{sec:deamsub} handles honestly.  Panel~(b) shows the premium
itself in dollars: at most \$\MacDeamWedgeMaxPremiumDollars, largest
deep in the money and fading through the forward (dashed line) toward
zero out of the money.  Hold the two panels side by side: a two-dollar
premium and a four-thousand-basis-point lie.  The quote is not wrong;
the model reading it is.

\begin{figure}[htbp]
\centering
\paperfig{\textwidth}{fig_deam_wedge.pdf}
\caption{The naive reading of American prices (synthetic calls, $6\%$
dividend yield against a $2\%$ rate, true volatility $25\%$, half a
year).  (a)~The same converged tree prices inverted two ways: the naive
European inversion books the premium as volatility --- up to
\MacDeamWedgeMaxBiasBp\ vol bp deep in the money --- while the
de-Americanization of \cref{sec:deamsub} recovers the flat truth (rms
\MacDeamWedgeRootRmsBp\ vol bp); below \$\MacDeamMapCallBdryPct\ the
price equals its exercise floor and no volatility exists to recover.  (b)~The premium
$A-E$ in dollars: at most \$\MacDeamWedgeMaxPremiumDollars.  Cents to
dollars on the left panel's axis become thousands of vol bp on the
right of panel (a) because the wings divide by a small vega.}
\label{fig:deamwedge}
\end{figure}

The chapter now walks the repair from first principles.
\Cref{sec:deamwhere} proves where the premium lives, so we know where
work is needed before doing any.  \Cref{sec:deamtree} builds the
pricing machine --- a binomial tree small enough to solve by hand.
\Cref{sec:deamsub} performs the subtraction and accounts for its
errors.  \Cref{sec:deamdivs} handles dividends, which enter as dates,
not rates.  \Cref{sec:deamloop} closes the loop with Chapter~6 ---
the forward needed the de-Americanization that needed the forward ---
and reproduces the residual bow.  \Cref{sec:deammatters} measures what
all of it is worth on the frozen gallery, on the population a fit
actually consumes and on the one it discards, and writes the
chapter's contract.
